\documentclass[11pt]{article}
\usepackage[margin=1in]{geometry}
\usepackage{amsmath,amssymb,amsthm}
\usepackage{graphicx}
\usepackage[hidelinks]{hyperref}

\newtheorem{theorem}{Theorem}
\newtheorem{corollary}[theorem]{Corollary}
\newtheorem{remark}{Remark}

\DeclareMathOperator{\se}{se}

\title{Bidual-Complemented Spaces With \texorpdfstring{$\se(X)=1$}{se(X)=1} Have the Self-Extension Property}
\author{FA Banach automatic math run}
\date{2026-07-18}

\begin{document}
\maketitle

\paragraph{Status.}
Partial result, likely valid. This proves a bidual-complemented subcase of an
open question from Garcia-Gutierrez, Garcia Pacheco, Piniella, and Rambla-Barreno
\cite{GGPPR}.

\paragraph{Source question.}
The source paper defines a Banach space $X$ to be $k$-self-extensible if, for
every subspace $Y\subseteq X$ and every bounded linear operator $T:Y\to Y$, one
can extend $T$ to an operator $S:X\to X$ satisfying
\[
  \|S\|\leq k\|T\|.
\]
It then sets
\[
  \se(X)=\inf\{k\geq 1: X\text{ is }k\text{-self-extensible}\}.
\]
After Corollary 5.4, the paper asks whether $X$ has the self-extension property
whenever $\se(X)=1$.

\begin{center}
\includegraphics[width=\textwidth]{figures/open_problem_crop.png}
\end{center}

\section*{Idea of the Proof}

The finite-dimensional proof in the source passes from almost isometric
extensions to an exact extension by compactness of bounded sets in
$\mathcal L(X,X)$. Infinite-dimensional operator balls are not norm compact, so
that argument does not directly survive.

For an arbitrary Banach space $X$, the canonical embedding $J_X:X\to X^{**}$
puts the values of almost optimal extensions inside weak-star compact balls of
$X^{**}$. For a fixed subspace $Y$ and operator $T:Y\to Y$, choose extensions
$S_n:X\to X$ with
\[
  \|S_n\|\leq (1+1/n)\|T\|.
\]
The family $(J_XS_nx)_{x\in X}$ has a pointwise weak-star cluster subnet in
$X^{**}$. The cluster point is a norm-preserving extension $R:X\to X^{**}$.
If there is a norm-one projection $P:X^{**}\to J_X(X)$, then $P\circ R$ lands
back in $X$ without increasing the norm. Thus the true obstruction to the
arbitrary case is not the compactness step, which always gives an $X^{**}$-valued
operator, but the lack of an exact contractive projection back to $X$.

\section*{Theorem}

\begin{theorem}
Let $X$ be a Banach space that is 1-complemented in its bidual; namely, suppose
there is a norm-one projection $P:X^{**}\to J_X(X)$, where
$J_X:X\to X^{**}$ is the canonical embedding. If $\se(X)=1$, then $X$ has the
self-extension property. Equivalently, for every subspace $Y\subseteq X$ and
every bounded linear operator $T:Y\to Y$, there is a bounded linear operator
$S:X\to X$ such that $S|_Y=T$ and $\|S\|=\|T\|$.
\end{theorem}

\begin{proof}
It is enough to consider closed subspaces. Indeed, if $Y$ is not closed and
$T:Y\to Y$ is bounded, then $T$ extends uniquely to a bounded operator
$\overline T:\overline Y\to \overline Y$ with the same norm, and any extension
of $\overline T$ restricts to an extension of $T$.

Fix a closed subspace $Y\subseteq X$ and $T\in\mathcal L(Y,Y)$. The case
$T=0$ is immediate, so assume $T\neq 0$. Since $\se(X)=1$, for each positive
integer $n$ there is an extension $S_n\in\mathcal L(X,X)$ of $T$ with
\[
  \|S_n\|\leq (1+1/n)\|T\|.
\]
Put $C=2\|T\|$. For each $x\in X$, the set
\[
  K_x=C\|x\| B_{X^{**}}
\]
is compact for the weak-star topology $\sigma(X^{**},X^*)$ by
Banach--Alaoglu. Hence the product
\[
  K=\prod_{x\in X} K_x
\]
is compact by Tychonoff's theorem. The points
\[
  \Phi_n=(J_XS_nx)_{x\in X}\in K
\]
therefore have a convergent subnet, say $\Phi_{n_\alpha}\to \Phi=(u_x)_{x\in X}$
in the product weak-star topology. Along such a subnet, $n_\alpha\to\infty$ in
the usual cofinal sense.

Define $R:X\to X^{**}$ by $Rx=u_x$. We first check linearity. Let $a,b$ be scalars
and $x,z\in X$. For every $x^*\in X^*$,
\[
\begin{aligned}
  u_{ax+bz}(x^*)
  &=\lim_\alpha x^*(S_{n_\alpha}(ax+bz)) \\
  &=a\lim_\alpha x^*(S_{n_\alpha}x)
    +b\lim_\alpha x^*(S_{n_\alpha}z) \\
  &=(aRx+bRz)(x^*).
\end{aligned}
\]
Since $X^*$ separates points of $X^{**}$, this gives
$R(ax+bz)=aRx+bRz$.

Next, $R$ extends $J_XT$. If $y\in Y$, then $S_ny=Ty$ for every $n$, hence
$Ry=J_XTy$.

Finally we prove the norm estimate. Fix $x\in X$ and $\varepsilon>0$. Since
$n_\alpha\to\infty$, eventually
\[
  \|J_XS_{n_\alpha}x\|\leq (1+\varepsilon)\|T\|\,\|x\|.
\]
The ball $(1+\varepsilon)\|T\|\,\|x\|B_{X^{**}}$ is weak-star closed, so the
weak-star limit $Rx$ belongs to it. Therefore
\[
  \|Rx\|\leq (1+\varepsilon)\|T\|\,\|x\|.
\]
Letting $\varepsilon\downarrow 0$ yields $\|Rx\|\leq \|T\|\,\|x\|$ for every
$x\in X$, so $\|R\|\leq\|T\|$.

Let $Q=J_X^{-1}P:X^{**}\to X$, where $J_X^{-1}$ is the inverse isometry from
$J_X(X)$ onto $X$. Define $S=Q\circ R$. Then
\[
  \|S\|\leq \|Q\|\|R\|\leq \|T\|.
\]
If $y\in Y$, then $Ry=J_XTy$ and $QJ_XTy=Ty$, so $S|_Y=T$. Hence also
$\|S\|\geq \|T\|$, and therefore $\|S\|=\|T\|$.
\end{proof}

\begin{corollary}
If $X$ is a dual Banach space and $\se(X)=1$, then $X$ has the self-extension
property.
\end{corollary}

\begin{proof}
Write $X=E^*$. Then $X^{**}=E^{***}$. The restriction map
\[
  E^{***}\longrightarrow E^*,\qquad \Phi\mapsto \Phi|_E,
\]
has norm one and restricts to the identity on the canonical copy of $E^*$ in
$E^{***}$. Thus $X$ is 1-complemented in $X^{**}$, and the theorem applies.
\end{proof}

\section*{Verification Notes and Limitations}

The proof uses only Banach--Alaoglu, Tychonoff compactness, weak-star
closedness of norm balls in $X^{**}$, and the assumed norm-one projection from
$X^{**}$ onto $X$. It does not require the extensions $S_n$ to be weak-star
continuous; it only takes weak-star limits of their point values after applying
the canonical embedding into $X^{**}$.

This is a partial answer to the source question. It covers all dual Banach
spaces, all reflexive spaces, and more generally every Banach space that is
contractively complemented in its bidual. It does not settle the same question
for arbitrary Banach spaces. It also does not address the two other source
questions about hyperplanes or spaces satisfying $\se(X)=\infty$.

\section*{Novelty Check}

The lightweight run indexes contained no exact entry for arXiv:1911.07826.
Local searches for the source title, self-extension property, self-extension
coefficient, $k$-self-extensible spaces, and $\se(X)=1$ found only the source
paper and unrelated occurrences. Web searches for the same terms, with
``dual'', ``bidual'', and ``reflexive'' variants, found the source paper
\cite{GGPPR} and the earlier Lyapunov/self-extension paper \cite{LPR}, but no
exact later answer to the bidual-complemented subcase. Novelty confidence is
therefore bounded rather than absolute.

\section*{Human Review Recommendation}

Check that the definition of subspace intended in the source permits the
initial closure reduction. If the source convention silently restricts to
closed subspaces, the reduction is unnecessary and the proof begins after it.
The main compactness step should be reviewed by verifying that the chosen
subnet has indices tending to infinity, so that the final norm of the
$X^{**}$-valued extension is exactly $\|T\|$, not merely $2\|T\|$. The final
projection step is exactly where the arbitrary Banach-space case remains open.

\begin{thebibliography}{9}

\bibitem{GGPPR}
J. B. Garcia-Gutierrez, F. J. Garcia Pacheco, P. Piniella, and
F. Rambla-Barreno,
\emph{Non-Hilbert Banach spaces with the self-extension property},
arXiv:1911.07826, 2019/2021.

\bibitem{LPR}
K. Loskot, A. Polanski, and R. Rudnicki,
\emph{Further comments on: Vector norms as Lyapunov functions for linear systems},
IEEE Transactions on Automatic Control 43 (1998), no. 2, 289--291.

\bibitem{Rudin}
W. Rudin,
\emph{Functional Analysis}, 2nd ed.,
McGraw--Hill, 1991.

\end{thebibliography}

\end{document}
